\documentclass{syllabus}

\begin{document}

\thispagestyle{fancy}
\date{Fall 2025}
\author{}
\title{PHY 4960: Portfolio Seminar D \\ {\large \mdseries Portfolio curation and reflection}}
\maketitle



\section*{General Information}

\raisebox{0.9\height}{

\begin{tabular}{lll}
\textbf{Instructor:}    & \multicolumn{2}{l}{Joseph Anderson}                             \\
\textbf{}               & \multicolumn{2}{l}{\href{mailto:janderson19@carthage.edu}{janderson19@carthage.edu}}                    \\
\textbf{}               & \multicolumn{2}{l}{DSC 082}                    \\
\textbf{}               & \multicolumn{2}{l}{(262)551-5871}                               \\ \\
\textbf{Office hours:}  & Wed, Fri               & 10:20 AM - 12:00 PM                    \\
\textbf{}               & Tuesday                & 3:30-5:00 PM (Discussion)              \\ \\
\textbf{Textbook:}      & \multicolumn{2}{l}{N/A} \\ \\
\textbf{Prerequisites:} & \multicolumn{2}{l}{PHY 3970 (Port. C)}                      
\end{tabular}
%
}

\section*{Course overview}

This is the fourth and final course in the Physics Portfolio Seminar sequence. This seminar focuses on 
preparing students to present and defend their physics portfolio, a collection of work that each student has 
produced while a physics major. Other elements of this seminar will include strategies for exploring post-graduate employment or further education, development of interview skills, and the production of 
effective resumes and CVs.

\subsection*{Student learning outcomes} 
Upon completion of this course, a student who engages thoughtfully and earnestly with the course activities will be able to:
\begin{enumerate}[label=\textcolor{CarthageDark}{\bfseries\arabic*.}]
    \item Describe and discuss the significance of their own work as a student toward their growth and development as an undergraduate physics major.
    \item Identify possible career or post-graduate pathways that align with or make use of the skills developed within the Carthage physics curriculum, broadly defined so as to include courses and co-curricular activities housed both within and outside of the Physics department.
    \item Identify and contextualize their own skills, aptitudes, and interests with respect to possible career paths or post-baccalaureate options.
    \item Use language appropriate for academic or professional environments.
\end{enumerate}



\section*{Course components}

\paragraph{Portfolio webpage} 
Students will create a portfolio showcasing personal qualities and skills as well as their growth in becoming a physicist. This will consist of 
\begin{itemize}
    \item a curated physics portfolio: body of \~5 significant works from the past 3 years.
    \item reflection pieces describing the physics and personal significance of key portfolio elements
    \item a context-rich inventory of skills, aptitudes, and interests
    \item a capstone reflection personal statement demonstrating a trajectory of growth and development as a Carthage physics major and identifying future career or post-graduate pathways
\end{itemize}

\paragraph{Portfolio defense}
The portfolio counts for the college thesis requirement; an oral presentation/defense of the webpage and reflection will be a major part of the grade in this course.

\paragraph{Professional materials}
Students will generate a resume, curriculum vitae (CV), and update/create their LinkedIn account profile.

\paragraph{Sample application}
Students will prepare a sample cover letter and sample application.

\paragraph{Other in-class assignments}
Other class activities will include: 
\begin{itemize}
    \item Written assignments
    \item Peer and group oral discussions
    \item Interviews, self-story, and elevator pitches
\end{itemize}

\paragraph{Interacting with PHY 1200} Interaction with PHY 1200 is encouraged. This may take the form of a) visiting the class to discuss club opportunities, b) help with instruction of the class or present a demonstration, c) helping out with a lab. Extra credit will be given for up to three 


\section*{Grading Scheme}
\subsection*{Component Weights}

\noindent\begin{tabular}{ll}
\textbf{Component} & \textbf{Weight}       \\ \hline
\textbf{Portfolio webpage}            & 40\%  \\
\textbf{Portfolio defense}            & 15\%       \\
\textbf{CV, etc.}            &     10\%   \\
\textbf{Sample application}            & 5\%       \\
\textbf{Participation and activities}            & 30\%       \\
\textbf{1200 interactions}            & +6\%       \\
\end{tabular}%



\lettergrades{a) authentic engagement with class discussions and reflection assignments, and b) pattern of seeking help in developing a professional development plan from the instructor, academic advisor, or other identified mentor.}

\section*{Policies}
\subsection{General course policies}
Expect a response within one business day for emails sent to the instructor. To ensure a prompt response, format the subject line as: [PHY 4960] *insert subject here*

Students are expected to arrive and leave on time. Cf. participation section for more information. I will attempt to arrive five minutes early and be available for at least five minutes after class.

\collegepolicies

\end{document}
